\subsection{1. Integer Sequences from Configurations in the Hausdorff Metric Geometry}
\textbf{OSTI:} \texttt{1997354}\quad\textbf{Authors:} Bobrowski, Elpers, Helmkamp, Ovsyannikov, Xique (2021)\quad\textbf{Domain:} Pure Mathematics / Combinatorics\\
\scorebadge{Coverage}{10}\quad\scorebadge{Agreement}{10}\quad\textbf{Total:} 20/20\par\smallskip
\noindent\textbf{Coverage rationale.} All 13 closed-form formulas from Theorems 8, 10, 11, 12 implemented and verified against brute-force bipartite-graph enumeration. All 21 OEIS sequences cross-referenced term-by-term. Exhaustive enumeration of 327,680 graphs up to K, and K, to confirm achievability of [1,100]. All five known non-achievable integers (19,37,41,59,67) confirmed. Fibonacci and Lucas achievability partially reconfirmed.\par
\noindent\textbf{Agreement rationale.} Every closed-form formula matches recursive enumeration exactly (integer arithmetic, no floating-point error). 21/21 OEIS sequences matched term-by-term. All 5 gaps confirmed. No discrepancies detected anywhere.\par\smallskip
\noindent\textbf{What reproduced:}
\begin{itemize}[leftmargin=1.4em,itemsep=0.2em,topsep=0.2em]
\item All 13 formula functions from Theorems 8, 10, 11, 12
\item Brute-force bipartite edge-cover enumeration up to K, and K,
\item OEIS cross-referencing for all 21 sequences (A048291, A335608--A335613, A337416--A337418, A340173--A340175, A340199--A340201, A340897--A340899, A342580, A342796)
\item Symmetry property E(m,n)=E(n,m)
\item Five known non-achievable integer gaps (19,37,41,59,67) in [1,100]
\item Fibonacci/Lucas achievability subsequences (8/10 and 5/6)
\end{itemize}\noindent\textbf{What missing:}
\begin{itemize}[leftmargin=1.4em,itemsep=0.2em,topsep=0.2em]
\item Extension beyond m,n$\leq$6 to validate higher-order formulas at scale
\item Complete resolution of the 35/100 unresolved integers in [1,100]
\item Proof-level verification (we only verify numerically, not symbolically)
\item Larger OEIS extensions to bound the next few non-achievable integers \textgreater{} 100
\item Asymptotic density analysis of achievable integers
\item Algorithmic characterization of non-achievability beyond brute force
\end{itemize}\noindent\textbf{Follow-on research questions:}
\begin{enumerate}[leftmargin=1.6em,itemsep=0.2em,topsep=0.2em,label=Q\arabic*.]
\item Extend the exhaustive graph enumeration to K, and K, to resolve the 35 open integers in [1,100] --- which remain non-achievable
\item Is there a polynomial-time certificate for non-achievability, or does deciding achievability require exponential search Empirically test with SAT encoding on integers up to 200.
\item Compute the asymptotic density of achievable integers in [1,N] as N$\rightarrow$$\infty$ --- does it tend to 1, to a constant \textless{} 1, or to 0
\item Generalize the bijection to tripartite graph edge covers --- what new integer sequences arise, and do they correspond to existing OEIS entries
\item Which formula family (from Theorems 10--12) is tightest --- i.e., attains the minimum m+n realization for a given integer Produce an atlas mapping each achievable integer to its cheapest bipartite realization.
\end{enumerate}

